\documentclass[11pt]{amsart}


\usepackage{amsmath, amssymb, amscd, cancel, graphicx, soul, stmaryrd, pstricks, pinlabel}
\usepackage{enumerate}
\usepackage{tikz-cd}

\usepackage{mathdots}

\headheight=7pt         \topmargin=14pt
\textheight=574pt       \textwidth=445pt
\oddsidemargin=18pt     \evensidemargin=18pt

\usepackage[all]{xy}
\usepackage{cleveref}

\setlength{\parskip}{4pt}


\newtheorem{thm}{Theorem}[section]
\newtheorem{conj}[thm]{Conjecture}
\newtheorem{rem}[thm]{Remark}
\newtheorem{cor}[thm]{Corollary}
\newtheorem{lem}[thm]{Lemma}
\newtheorem{prop}[thm]{Proposition}
\newtheorem{ques}[thm]{Question}
\newtheorem*{ques*}{Question}
\newtheorem{defin}[thm]{Definition}
\newtheorem{construction}[thm]{Construction}
\newtheorem{example}[thm]{Example}
\newtheorem*{example*}{Example}

\numberwithin{equation}{section}

\clubpenalty=3000
\widowpenalty=3000

\def\cB{{\mathcal B}}
\def\cC{{\mathcal C}}
\def\cF{{\mathcal F}}
\def\cH{{\mathcal H}}
\def\q{{\mathcal Q}}
\def\s{{\mathfrak s}}
\def\t{{\mathfrak t}}
\def\x{{{\bf x}}}
\def\y{{{\bf y}}}

\def\bC{{\mathbb C}}
\def\D{{\mathcal D}}
\def\E{{\mathcal E}}
\def\bF{{\mathbb F}}
\def\G{{\mathcal G}}
\def\H{{\mathbb H}}
\def\M{{\mathcal M}}
\def\MP{{\mathcal M}_{\rm P}}
\def\MR{{\mathcal M}_{\rm R}}
\def\MD{{\mathcal M}_{\rm D}}
\def\barMD{\overline{{\mathcal M}_{\rm D}}}
\def\barMP{\overline{{\mathcal M}_{\rm P}}}
\def\barMR{\overline{{\mathcal M}_{\rm R}}}

\def\Cut{{\mathrm{Cut}}}
\def\Flow{{\mathrm{Flow}}}

\def\O{{\mathcal O}}
\def\P{{\mathcal P}}
\def\bQ{{\mathbb Q}}
\def\bR{{\mathbb R}}
\def\T{{\mathcal T}}
\def\bZ{{\mathbb Z}}

\def\L{{\langle}}
\def\R{{\rangle}}

\def\Char{{\mathrm{Char}}}
\def\Cut{{\mathrm{Cut}}}
\def\Div{{\mathrm{Div}}}
\def\Flow{{\mathrm{Flow}}}
\def\Hom{{\mathrm{Hom}}}
\def\Irr{{\text{Irr}}}
\def\Pic{{\text{Pic}}}
\def\spc{{\mathrm{spin^c}}}
\def\Spc{{\mathrm{Spin^c}}}
\def\Short{{\mathrm{Short}}}
\def\Sym{{\mathrm{Sym}}}

\DeclareMathOperator{\disc}{disc}
\DeclareMathOperator{\im}{im}
\def\cok{{\mathrm{coker}}}
\def\cut{{\textup{cut}}}
\def\cyc{{\textup{cyc}}}
\def\del{{\partial}}
\def\mod{{\textup{mod} \;}}
\def\rk{{\mathrm{rk}}}

\def\supp{{\textup{supp}}}
\def\Vor{{\textup{Vor}}}

\newcommand{\into}{\hookrightarrow}
\newcommand{\onto}{\twoheadrightarrow}
\newcommand{\ol}{\overline}
\newcommand{\norm}[1]{\left| #1 \right|^2}

\title[Problem Submission]{First Proof $\bullet$ Second Batch \\ Problem Submission}

\author{Joshua Evan Greene}
\address{Boston College, USA}
\email{joshua.greene@bc.edu}
\author{Duncan McCoy}
\address{Universit\'e du Qu\'ebec \`a Montr\'eal, Canada}
\email{mc\_coy.duncan@uqam.ca}

\begin{document}
	
\maketitle

\section*{Problem Statement.}

Let $T$ denote a finite, weighted tree with vertex set $V$, edge set $E$, and weight function $w : V \longrightarrow \bZ$.
Let $L(T)$ denote the free abelian group generated by $V$, equipped with the symmetric, bilinear form $\cdot$ whose value on a pair of generators $u,v \in V$ is given by 
\[
u \cdot v = 
\begin{cases}
    w(u), & u=v; \\
    -1, & (u,v) \in E; \textup{ and} \\
    0, & u \ne v, (u,v) \notin E.
\end{cases}
\]
Suppose that the form is positive definite, and suppose that there exists a single vertex $v$ for which the vertex weight $w(v)$ is less than the vertex degree $d(v)$.
Prove that $T$ contains a vertex $u$ such that $u$ is an irreducible element of $L(T)$, i.e.~there do not exist nonzero lattice elements $a,b \in L(T)$ such that $u = a+b$ and $a \cdot b \ge 0$.


\section{Context}

A finitely generated, free abelian group equipped with a real-valued, symmetric bilinear form is called a {\em lattice}.
It is {\em integral} if the form takes integer values, {\em rational} if the form takes rational values, and {\em positive definite} if the form is positive definite.
An element of a positive definite lattice $L$ is called {\em reducible} if it is possible to express it as a sum of nonzero elements $a,b \in L$ such that $a \cdot b \ge 0$, and it is called {\em irreducible} otherwise.
For example, an element $a \in L$ of minimal nonzero self-pairing $|a|^2 = a \cdot a$ in a lattice $L$ is irreducible, such as an element of self-pairing 1 in an integral lattice, and the irreducible elements in the lattice of integer points in Euclidean space are precisely the elements of self-pairing 1.
Irreducible lattice elements specify the facets of the Voronoi polyhedron of $L$.
The lattices $L(T)$ considered in the problem statement are called {\em tree lattices}.
A vertex $v$ in a weighted tree is {\em bad} if $w(v) < d(v)$.
In this language, the problem statement is to prove:

\begin{thm}
\label{thm:one_bad_vtx}
If a weighted tree $T$ contains a single bad vertex and $L(T)$ is positive definite, then $L(T)$ contains an irreducible vertex element.
\end{thm}

\noindent
Theorem~\ref{thm:one_bad_vtx} is motivated by a conjecture of Walter Neumann and Don Zagier \cite{NeumannZagier1985,Neumann1989}:
\begin{conj}
\label{conj: nz}
If $T$ is a weighted tree and $L(T)$ is isometric to the lattice of integer points in Euclidean space, then there exists a vertex $v \in V$ with weight $w(v) = 1$.
\end{conj}

Conjecture~\ref{conj: nz} stems from a problem in low-dimensional topology concerning integer homology cobordism and gauge theory.
The original version of Conjecture~\ref{conj: nz} in \cite{NeumannZagier1985} is posed as a question for star-shaped trees, i.e.~those containing a single vertex $v$ with $d(v) \ge 3$.
Issa and McCoy settled this conjecture in \cite{IssaMcCoy2020}.
Neumann posed the more general Conjecture~\ref{conj: nz} in \cite{Neumann1989}, where he attributed it to himself and Zagier.

Generalizing \cite{IssaMcCoy2020}, we proved Conjecture~\ref{conj: nz} under the assumption that $T$ has a single bad vertex (compare \cite[Definition 1.1]{OzsvathSzabo2003}).
We then realized that we could prove the more general Theorem~\ref{thm:one_bad_vtx}.
It leads us to conjecture the following improvement on Conjecture~\ref{conj: nz}:

\begin{conj}
\label{conj: gm}
If $T$ is a weighted tree and $L(T)$ is positive definite, then at least one vertex of $T$ is irreducible in $L(T)$.
\end{conj}


We have two primary motivations for proposing this problem. Firstly, we believe that this problem has several characteristics making it a good AI test case. Although Theorem~\ref{thm:one_bad_vtx} can be proven with relatively little background, our proof involves several subtle steps.
Furthermore, there is very little literature relating to Conjecture~\ref{conj: nz}, and our proof of Theorem~\ref{thm:one_bad_vtx} does not resemble the techniques used in \cite{IssaMcCoy2020}.
Secondly, we are simply curious to see whether an AI attempt could shed fresh light on Conjectures~\ref{conj: nz} and~\ref{conj: gm}.

\section{Preparation and lemmas.}

\subsection{Cut and flow lattices}
We review some basic definitions about cut and flow lattices on graphs \cite{GodsilRoyle2001, Greene2013}.
Let $G$ be a finite, loopless graph with vertex set $V=V(G)$ and edge set $E=E(G)$, possibly with parallel edges.
We use $e(v,w)$ to denote the number of edges between a pair of distinct vertices $v$ and $w$ in $G$.

Choosing an arbitrary orientation for the edges of $G$ endows $G$ with the structure of a CW-complex. This allows us to consider the integral cellular chain complex
\begin{equation}\label{eq:cellular_chain_complex}
0\longrightarrow C_1(G) \overset{\partial}{\longrightarrow} C_0(G)\longrightarrow 0,
\end{equation}
where $\partial (e)= v-w$ for an edge $e$ oriented from $w$ to $v$.
We equip $C_1(G)$ and $C_0(G)$ with bilinear pairings such that $E$ and $V$ are orthonormal bases, respectively. These pairings give rise to the adjoint map
\begin{align}\begin{split}   
\partial^* : C_0(G)&\longrightarrow C_1(G),\\
v&\longmapsto \sum_{e \in E} (\partial(e) \cdot v) e.
\end{split}
\end{align}

The \emph{cut lattice} on $G$ is defined to be 
\[
\Cut(G)=\im \partial^* \subseteq C_1(G).
\]
Clearly the images of the vertices span $\Cut(G)$ and for any pair of vertices, $u,v\in V$, we have
\begin{equation}\label{eq:coboundary_graph_embedding}
\partial^* u \cdot \partial^* v = \begin{cases}
    d(u) &\text{if $u=v$}\\
    -e(u,v)&\text{if $u\neq v$}.
\end{cases}
\end{equation}
The \emph{flow lattice} of $G$ is defined to be
\[
\Flow(G)=\ker \partial \subseteq C_1(G).
\]
Since $\partial^*$ is adjoint to $\partial$, the lattices $\Cut(G)$ and $\Flow(G)$ are orthogonal to one another in $C_1(G)$.

If $L$ is an integral lattice, then $L \otimes \mathbb{Q}$ is a vector space with an induced symmetric, bilinear form.
Within it sits the dual lattice $L^* = \{ \lambda \in L \otimes \mathbb{Q} : \lambda \cdot v \in \bZ, \forall v \in L \}$.
A sublattice $L' \subset L$ is {\em primitive} if it satisfies any of the equivalent conditions:
\begin{itemize}
    \item 
    $L \cap (L' \otimes \bQ) = L'$, the intersection taking place within $L \otimes \bQ$;
    \item
    $L'$ admits a basis which completes to one of $L$; and
    \item
    the restriction/projection map $L^* \to (L')^*$ surjects.
\end{itemize}
The lattice $C_1(G)$ is self-dual, and $\Cut(G)$ and $\Flow(G)$ are both primitive sublattices of $C_1(G)$ \cite[Proposition 3.1]{Greene2013}.
The inclusion of the orthogonal direct sum $\Cut(G) \oplus \Flow(G) \hookrightarrow C_1(G)$ induces an inclusion of dual lattices $C_1(G) \hookrightarrow \Cut(G)^* \oplus \Flow(G)^*$.
We uniquely express every $x \in C_1(G)$ in the form $x = x_C + x_F$ with respect to this decomposition.

\subsection{The projection length formula}

The result of this section is a formula for the self-pairing of the projection $x_C \in \Cut(G)^*$ of an element $x \in C_1(G)$.

Let $G$ be a connected graph with a choice of orientation on its edges.
Let $(A,B)$ be an ordered partition of the vertex set of $G$. The cut vector defined by $(A,B)$ is
\[
\chi_{A,B}=\sum_{e\in E(A,B)} \varepsilon_{A,B,e} e \in C_1(G),
\]
where 
\[
\varepsilon_{A,B,e}=\begin{cases} 1 &\text{if $e$ oriented from $A$ to $B$}\\
 -1 &\text{if $e$ oriented from $B$ to $A$}.
\end{cases}
\]
Evidently, we have that $\chi_{A,B}=-\chi_{B,A}$. Observe also that $\chi_{A,B}$ can be expressed as
\begin{equation}\label{eq:chi_formula}
    \chi_{A,B}=-\sum_{v\in A} \partial^* v.
\end{equation}
This implies that $\chi_{A,B}\in \Cut(G)$.

A two-component spanning forest of $G$ is an unordered pair $\{A,B\}$ of subgraphs of $G$ such that $A$ and $B$ are trees and $A$ and $B$ partition the vertices of $G$.
 Given an element $x\in C_1(G;\bR)$, we define its square-flux on $\{A,B\}$ to be 
\[\mathcal{F}_{\{A,B\}}(x)=|\chi_{A,B}\cdot x|^2 = |\chi_{B,A}\cdot x|^2.\]
Explicitly, if $x$ takes the form $x=\sum_e x_e e$, this is
\begin{equation}\label{eq:flux_formula}
\mathcal{F}_{\{A,B\}}(x)=\sum_{e\in E(A,B)} x_e^2 + \sum_{\substack{e,e'\in E(A,B)\\ e\neq e'}} 2\varepsilon_{A,B,e}\varepsilon_{A,B,e'} x_e x_{e'}.
\end{equation}

\begin{lem}
\label{lem:length}
Let $G$ be a connected graph and let $x\in C_1(G)$. If $x=x_C+x_F$, where $x_C\in \Cut(G)^*$ and $x_F\in \Flow(G)^*$, then
\[
\norm{x_C}= \frac{1}{\tau} \sum_{\{A,B\}} \mathcal{F}_{\{A,B\}}(x),
\]
where the sum is over all two-component spanning forests and $\tau$ is the number of spanning trees in $G$.
\end{lem}

\noindent
{\em Remark.}
One can deduce Lemma~\ref{lem:length} from an expression for $x_C$ due to Kirchhoff \cite{Kirchhoff1847}.
(For a modern treatment, see \cite[Theorem 14.8.1]{GodsilRoyle2001}.)
We give a self-contained argument in the spirit of ``completing the square."

\begin{proof}
A \emph{connected spanning unicyclic subgraph} $R$ of $G$ is a subgraph which is connected, which contains every vertex of $G$, and which contains a unique cycle $C_R$.
Deleting any edge of this cycle therefore results in a spanning tree for $G$. Given a connected spanning unicyclic subgraph $R$, we define a circulation vector for $R$ to be a vector of the form $\omega_R = \sum_{e\in C_R} \delta_{R,e} e$, where the $\delta_{R,e}\in \{\pm 1\}$ are chosen so that $\partial \omega_R=0$. Note that the circulation vector is unique up to an overall change of sign. Given an element $x\in C_1(G;\bR)$, we define its circulation on $R$ to be 
\[\mathcal{C}_{R}(x)=|\omega_R\cdot x|^2.\]
The circulation can be calculated by the formula
\begin{equation}\label{eq:circulation_formula}
\mathcal{C}_{R}(x)=\sum_{e\in C_R} x_e^2 + \sum_{\substack{e,e'\in C_R\\ e\neq e'}} 2\delta_{R,e}\delta_{R,e'} x_e x_{e'}.
\end{equation}
We notice a number of natural bijections:
\begin{itemize}
    \item Given a fixed edge $e$, the set of two-component spanning forests $\{A,B\}$ with $e\in E(A,B)$ is in bijection with the set of spanning trees containing $e$.
    \item The set of connected spanning unicyclic subgraphs $R$ with $e\in C_R$ is in bijection with the set of spanning trees not containing $e$.
    \item Given a pair of distinct edges $e$ and $e'$, the set of two-component spanning forests $\{A,B\}$ with $e,e'\in E(A,B)$ is in bijection with connected spanning unicyclic subgraphs $R$ such that $e,e'\in C_R$. We obtain $\{A,B\}$ by deleting $e$ and $e'$ from $R$. Moreover note that under this bijection we always have $\varepsilon_{A,B,e}\varepsilon_{A,B,e'}=-\delta_{R,e}\delta_{R,e'}$.
\end{itemize}
 Thus summing \eqref{eq:flux_formula} and \eqref{eq:circulation_formula} over all two-component spanning forests and all connected spanning unicyclic subgraphs we obtain
\begin{equation}\label{eq:full_length}
\tau \norm{x}=\sum_{\{A,B\}}\mathcal{F}_{\{A,B\}}(x) + \sum_{R}\mathcal{C}_{R}(x).
\end{equation}
Since $x_C$ is orthogonal to every element of $\Flow(G)$, we have $\cC_R(x_C)=0$ for every connected spanning unicyclic subgraph $R$. Thus using \eqref{eq:full_length} to calculate $\norm{x_C}$ we obtain
\begin{equation}\label{eq:full_length2}
\tau \norm{x_C}=\sum_{\{A,B\}}\mathcal{F}_{\{A,B\}}(x_C).
\end{equation}
On the other hand, $x_F$ is orthogonal to every element of $\Cut(G)$. Consequently, we have
\[
\cF_{\{A,B\}}(x)=\cF_{\{A,B\}}(x_C)
\]
for any two-component spanning forests $\{A,B\}$, allowing us to obtain the desired formula from \eqref{eq:full_length2}.
\end{proof}


\subsection{Apex trees}

If $T$ is a weighted tree, then its vertex set $V$ is a basis for $L(T)$.
Let $V^*$ denote the dual basis for $L(T)^*$, and let $v^* \in V^*$ denote the dual basis element dual to $v \in V$.

\begin{lem}\label{lem:apex_tree}
    Let $T$ be a weighted tree with no bad vertices. Let $v$ be a vertex such that $w(v)>d(v)$. If $a\in L(T)^*$ satisfies $a\cdot v\neq 0$, then $\norm{a} \geq \norm{v^*}$.
\end{lem}

\begin{proof}
    Let $G$ be the graph obtained by adding to $T$ a single vertex $v_\infty$ joined by $w(u) - d(u) \ge 0$ edges to each vertex $u\in T$.
    We call such a graph $G$ an {\em apex tree} and $v_\infty$ its {\em apex}.
    By construction, $L(T)$ is isomorphic to the cut lattice $\Cut(G)$ via $v \longleftrightarrow \del^*(v)$.
    By hypothesis, there is at least one edge $e$ between $v$ and $v_\infty$.

    The projection of $e \in C_1(G)$ to $\Cut(G)^* = L(T)^*$ is the dual element $v^*$.
    Its square-flux across a two-component spanning forest $\{A,B\}$ of $G$ is 1, if  $A \cup B \cup e$ is a spanning tree of $G$, and 0, if not.
    By the projection length formula Lemma~\ref{lem:length}, it follows that
    \begin{equation}
        \label{e:self-pairing}
    \norm{v^*}=\tau_e/\tau
    \end{equation}
    where $\tau_e$ counts the number of spanning trees for $G$ containing $e$.

    Select any $a \in L(T)^*$ such that $a \cdot v \ne 0$.
    Write $x_C \in \Cut(G)^*$ for its image under the identification $\Cut(G)^* = L(T)^*$.
    Because the restriction map from $C_1(G)$ to $\Cut(G)^*$ surjects, there exists $x\in C_1(G)$ such that $x = x_C + x_F \in \Cut(G)^* \oplus \Flow(G)^*$.
    We aim to show that there are at least $\tau_e$ two-component spanning forests such that $\cF_{\{A,B\}}(x)\neq 0$.
    Then another application of Lemma~\ref{lem:length} yields $\norm{a} \ge \norm{v^*}$.
    
    To achieve our aim, we construct an injection $\phi$ from spanning trees containing $e$ to two-component spanning forests $\{A,B\}$ for which $\cF_{\{A,B\}}(x)\neq 0$.
    Let $R$ be a spanning tree containing $e$.
    Let $R_1, \dots, R_k$ be the connected components of $R-v$. We will consider the two-component spanning forests of the form $\{R_j,R-R_j\}$. These are precisely the spanning forests obtained from $R$ by deleting an edge incident to $v$. Since every vertex of $G-v$ is contained in precisely one $R_j$, \eqref{eq:chi_formula} implies that we have the following linear relation among cut vectors:
    \[
    \chi_{G-v,v}=\sum_{j=1}^k\chi_{R_j,R-R_j}.
    \]
    Since $\chi_{G-v,v}\cdot x=a\cdot v\neq 0$, this implies that there exists an index $\ell$ such that $\chi_{R_\ell,R-R_\ell}\cdot x\neq 0$. We define $\phi(R)=\{R_\ell,R-R_\ell\}$ for such an $\ell$ and observe that this satisfies the desired property that $\cF_{\{R_\ell,R-R_\ell\}}(x)$ is nonzero. It remains to verify that $\phi$ is injective. We will show how to reconstruct $R$ from $\phi(R)$, using the fact that $\phi(R)$ is a two-component spanning forest obtained from $R$ by deleting an edge incident to $v$. Let $\phi(R)=\{A,B\}$. If neither $A$ nor $B$ contains the edge $e$, then $R$, being a tree containing $e$, must take the form $R=A\cup B\cup e$. Otherwise w.l.o.g.~ we assume that $e$ is contained in $A$. This implies that $v, v_\infty\in A$ and hence that $B$ is a subtree of $T$ not containing $v$. Moreover, there is a unique edge $f$ between $v$ and a vertex of $B$, since $T$ is a tree. In this case $R$ must take the form $R=A\cup B\cup f$. Thus we have verified that $\phi$ is an injection. 
\end{proof}

\begin{cor}
\label{cor:apex_tree}
Under the assumptions of Lemma~\ref{lem:apex_tree}, the element $v^* \in L(T)^*$ is irreducible.
\end{cor}

\begin{proof}
    Suppose that $v^* = a+b$ with $a,b \in L(T)^*$ both nonzero.
    Then $1 = v^* \cdot v = (a+b) \cdot v = a \cdot v + b \cdot v$.
    Hence we may assume w.l.o.g.~that $a \cdot v \ne 0$.
    We have $\norm{a}\ge \norm{v^*}$ by Lemma~\ref{lem:apex_tree}, while $\norm{b} > 0$, since $b \ne 0$.
    Then
    \[
    \norm{v^*}= \norm{a+b}= \norm{a^*}+ 2(a \cdot b) + \norm{b^*}> \norm{v^*}+ 2 (a \cdot b) + 0.
    \]
    Therefore, $a \cdot b < 0$.
    This proves that $v^*$ is irreducible.
\end{proof}

\begin{cor}
\label{cor:tight}
Under the assumptions of Lemma~\ref{lem:apex_tree}, $0 < |v^*|^2 \le 1$, and if $|v^*|^2 =1$, then either $w(v) = 1$ or $T$ contains a leaf $u \ne v$ with $w(u)=1$.
\end{cor}

\begin{proof}
    Equation \eqref{e:self-pairing} shows that $0 < |v^*|^2 \le 1$, and $|v^*|^2=1 \iff \tau_e = \tau \iff$ every spanning tree of $G$ uses $e \iff e$ is a cut-edge of $G \iff  e$ is the sole edge of $G$ incident with the apex vertex $\iff$ $w(v) = d(v)+1$ and $w(u) = d (u)$ for all $u \in V$, $u \ne v$.
    If $V = \{v\}$, then $w(v)=1$; if not, then $T$ contains a leaf $u \ne v$, and $w(u) = d(u) = 1$.
\end{proof}

\subsection{Reducible vertices.}
\label{sec:reducible}
Let $T$ be a weighted tree such that $L(T)$ is positive definite. Let $v$ be a vertex in $T$, and let $T_1,\dots, T_k$ be the weighted trees obtained by deleting $v$ from $T$.
For each tree $T_i$, let $v_i\in T_i$ be the vertex adjacent to $v$.
Let $\lambda_i \in L(T_i)^*$ denote the dual to $v_i$ with respect to the vertex basis $V(T_i)$.
We regard $\lambda_i \in L(T) \otimes \bQ$ under the inclusion $L(T_i)^* \subset L(T_i) \otimes \bQ \subset L(T) \otimes \bQ$.
Equivalently, $\lambda_i \in L(T) \otimes \bQ$ is the orthogonal projection to $L(T_i) \otimes \bQ$ of the element $v_i^* \in L(T)^*$ dual to $v_i$ with respect to the vertex basis $V(T)$.

Form the sublattice $L_1 = L(T_1) \oplus \cdots \oplus L(T_k) \subset L(T)$. 
It is primitive, because the basis $V(T) \setminus \{v\}$ for it completes to the basis $V(T)$ of $L(T)$.
Its orthogonal complement $L_1^\perp \subset L(T) \otimes \bQ$ is a 1-dimensional rational vector space with a symmetric, bilinear form.
Let $\pi : L(T) \rightarrow L_1^\perp$ denote the orthogonal projection, and let $L_0=\pi(L(T))$ be the image of the projection.
It is a rational lattice.
Thus, $L(T)$ is contained in the orthogonal direct sum $L_0 \oplus L_1^* = L_0 \oplus L(T_1)^* \oplus \cdots \oplus L(T_k)^*$.

\begin{lem}
\label{lem:irred}
    The projection of $v$ to $L_0$ is irreducible in $L_0$.
\end{lem}

\begin{proof}
Every vertex of $T$ besides $v$ projects to 0 in $L_0$.
Therefore, $\pi(v)$ generates $L_0$, hence is irreducible.
\end{proof}

\begin{lem}
\label{lem:reducible_dichotomy}
    If $v$ is reducible, then there exists an index $1\leq i\leq k$ such that either
    \begin{enumerate}[(i)]
        \item $\lambda_i$ is reducible in $L(T_i)^*$, or
        \item $\lambda_i\in L(T_i)$.
    \end{enumerate}
\end{lem}

\begin{proof}
Because $v \cdot v_i = -1$ for all $i$ and $v \cdot u = 0$ for every $u \in V(T_i)$, $u \ne v_i$, it follows that the projection of $v$ to $L(T_i)^*$ is $-\lambda_i$ for all $i$.
We therefore obtain the orthogonal decomposition
\[
v=\pi(v)-\lambda_1-\dots- \lambda_k.
\]
If $v$ is reducible, then we can decompose $v=x+y$ with nonzero $x,y\in L(T)$ satisfying $x\cdot y\geq 0$.
These have orthogonal decompositions
\[
\text{$x=\overline{x}+ x_1 +\dots + x_k$ \quad and \quad $y=\overline{y}+ y_1 +\dots + y_k$}
\]
with $\overline{x},\overline{y}\in L_0$ and $x_i,y_i\in L(T_i)^*$ for $i=1,\dots, k$.
Thus, $\overline{x}+\overline{y}=\pi(v)$ and $x_i+y_i=-\lambda_i$ for all $i$.
Since $\pi(v)$ is irreducible in $L_0$ by Lemma~\ref{lem:irred}, we have  $\overline{x}\cdot \overline{y}\leq 0$, with equality only if $\overline{x}=0$ or $\overline{y}=0$.
On the other hand,
\[
x\cdot y= \overline{x}\cdot \overline{y} + \sum_{i=1}^k x_i\cdot y_i\geq 0.
\]

If $\overline{x}\cdot \overline{y}<0$, then there exists an index $i$ such that $x_i\cdot y_i>0$, so $\lambda_i$ is reducible in $L(T_i)^*$, which returns conclusion (i).
If instead no $\lambda_i$ is reducible in $L(T_i)^*$, then $\overline{x} \cdot \overline{y} = 0$ and $x_i \cdot y_i = 0$ for all $i$.
Hence $\{ \overline{x}, \overline{y} \} = \{ 0 , \pi(v) \}$ and $\{ x_i, y_i \} = \{ 0 , -\lambda_i \}$ for all $i$.
Assume w.l.o.g.~that $\overline{x}=0$.
Then $x$ takes the form
\[
x=-\sum_{i\in A} \lambda_i
\]
for some $A\subseteq \{1,\dots, k\}$, and $A \ne \emptyset$, since $x \ne 0$.
Because $L_1 \subset L(T)$ is primitive and $x \in L(T) \cap (L_1 \otimes 
\bQ)$, it follows that $x \in L_1$.
Because the $\lambda_i$ are drawn from orthogonal summands of $L_1^*$, it follows that $\lambda_i \in L(T_i)$ for all $i \in A \ne \emptyset$, which returns conclusion (ii).
\end{proof}

\section{Final\'e}

\begin{proof}[Proof of Theorem~\ref{thm:one_bad_vtx}]
 We argue somewhat more: either the bad vertex $v$ is irreducible, or else $T$ contains a leaf of weight 1 (which is therefore irreducible).

Define trees $T_i$ and vertices $v_i$ as in Section \ref{sec:reducible}.
The tree $T_i$ has no bad vertices by assumption.
In addition, the degree of $v_i$ in $T_i$ is one less than its degree in $T$, which is $\le w(v_i)$.
Therefore, Corollary~\ref{cor:apex_tree} implies that $\lambda_i \in L(T_i)^*$ is irreducible for all $i$.

If $\lambda_i \notin L(T_i)$ for all $i$, then Lemma~\ref{lem:reducible_dichotomy} implies that $v$ is irreducible, as desired.

If instead $\lambda_i \in L(T_i)$ for some $i$, then $|\lambda_i|^2 \in \bZ$.
By Corollary \ref{cor:tight}, we must have $|\lambda_i|^2=1$, and either $w(v_i) = 1$ or $T_i$ contains a leaf $u \ne v_i$ with $w(u)=1$.
Either case yields a leaf of $T$ of weight 1, as desired.
\end{proof}

\noindent
{\em Example:} Let $T$ be the weighted tree with a single vertex $v$ with $d(v) = 3 $ and $w(v) = 2$ and three leaves of weights 1, 2, and 3.  Then $T$ contains a single bad vertex $v$, and $L(T)$ is positive definite (it is isometric to the lattice of integer points in Euclidean space), but $v$ is reducible.


\providecommand{\bysame}{\leavevmode\hbox to3em{\hrulefill}\thinspace}
\providecommand{\MR}{\relax\ifhmode\unskip\space\fi MR }
\providecommand{\MRhref}[2]{%
  \href{http://www.ams.org/mathscinet-getitem?mr=#1}{#2}
}
\providecommand{\href}[2]{#2}
\begin{thebibliography}{Neu89}

\bibitem[GR01]{GodsilRoyle2001}
Chris Godsil and Gordon~F. Royle, \emph{Algebraic graph theory}, Graduate Texts
  in Mathematics, vol. 207, Springer, New York, 2001.

\bibitem[Gre13]{Greene2013}
Joshua~E. Greene, \emph{Lattices, graphs, and {C}onway mutation}, Inventiones
  Mathematicae \textbf{192} (2013), no.~3, 717--750.

\bibitem[IM20]{IssaMcCoy2020}
Ahmad Issa and Duncan McCoy, \emph{On {S}eifert fibered spaces bounding
  definite manifolds}, Pacific Journal of Mathematics \textbf{304} (2020),
  no.~2, 463--480.

\bibitem[Kir47]{Kirchhoff1847}
Gustav Kirchhoff, \emph{{Ueber die Aufl{\"o}sung der Gleichungen, auf welche
  man bei der Untersuchung der linearen Vertheilung galvanischer Str{\"o}me
  gef{\"u}hrt wird}}, Annalen der Physik und Chemie \textbf{72} (1847),
  497--508.

\bibitem[Neu89]{Neumann1989}
Walter~D. Neumann, \emph{On bilinear forms represented by trees}, Bulletin of
  the Australian Mathematical Society \textbf{40} (1989), no.~2, 303--321.

\bibitem[NZ85]{NeumannZagier1985}
Walter~D. Neumann and Don Zagier, \emph{A note on an invariant of {F}intushel
  and {S}tern}, Geometry and Topology, Lecture Notes in Mathematics, vol. 1167,
  Springer, Berlin, 1985, Proceedings of the Special Year in Topology,
  University of Maryland, College Park, 1983/84, pp.~241--244.

\bibitem[OS03]{OzsvathSzabo2003}
Peter Ozsv{\'a}th and Zolt{\'a}n Szab{\'o}, \emph{On the {F}loer homology of
  plumbed three-manifolds}, Geometry \& Topology \textbf{7} (2003), 185--224.

\end{thebibliography}

\end{document}